\section{Aritmética de vacancias entre las escalas \texorpdfstring{\(3^6\)}{3 elevado a 6} y \texorpdfstring{\(10^3\)}{10 elevado a 3}}
\label{sec:vacancias-corriente-electron-rev11}
\HMTClaim{MD-R-VACANCY-POLARIZATION-CURRENT}
\index{vacancia!calendario aritmético}
\index{polarización!corriente discreta}

Las escalas finitas
\[
 729=3^6,
 \qquad
 1000=10^3
\]
determinan un desajuste logarítmico exacto. Definimos
\begin{equation}
 \rho_{\rm vac}=\log_{1000}729=\log_{10}9,
 \qquad
 \delta_{\rm vac}=1-\rho_{\rm vac}=\log_{10}(10/9).
 \label{eq:delta-vacancias-rev11}
\end{equation}
La irracionalidad de \(\delta_{\rm vac}\) se obtiene directamente de las
valuaciones primas: una igualdad
\(\delta_{\rm vac}=p/q\) implicaría
\((10/9)^q=10^p\), incompatible con las valuaciones en \(2,3,5\).
La vacancia es, por tanto, una rotación aritmética aperiódica determinada
íntegramente por la relación entre las dos escalas; su definición precede a
la ley electrónica.

\subsection{Leyes exactas de hueco y retorno}

Sea
\[
 p_n=\left\lfloor\frac{n}{\delta_{\rm vac}}\right\rfloor.
\]

\begin{theorem}[Ley de huecos \texorpdfstring{\(21/22\)}{21/22}]
\label{thm:ley-huecos-vacancias-rev11}
Los incrementos consecutivos de la sucesión \(p_n\) satisfacen
\[
 p_{n+1}-p_n\in\{21,22\}.
\]
Ambos valores aparecen en la palabra mecánica.
\end{theorem}

\begin{proof}
Para todo irracional \(x>1\), la diferencia
\(\lfloor(n+1)x\rfloor-\lfloor nx\rfloor\) pertenece a
\(\{\lfloor x\rfloor,\lceil x\rceil\}\). En este caso
\[
 \delta_{\rm vac}^{-1}
 =21.85434532678283256\ldots,
\]
lo que determina los dos huecos. La irracionalidad excluye la estabilización
en uno solo de ellos.
\end{proof}

La frecuencia del hueco corto es
\[
 \sigma_{\rm vac}=22-\delta_{\rm vac}^{-1}
 =0.1456546732171674374\ldots,
 \qquad
 \sigma_{\rm vac}^{-1}=6.865553832996660997\ldots.
\]

\begin{corollary}[Ley de retorno \texorpdfstring{\(6/7\)}{6/7}]
\label{cor:ley-retorno-vacancias-rev11}
Las apariciones sucesivas del hueco \(21\) están separadas por seis o por
siete posiciones.
\end{corollary}

\begin{proof}
La palabra indicadora del hueco corto es la codificación mecánica de
pendiente irracional \(\sigma_{\rm vac}\). El argumento de Beatty aplicado a
\(\sigma_{\rm vac}^{-1}\) produce exactamente los retornos \(6/7\)
\cite{Lothaire2002}.
\end{proof}

Las dos leyes forman el marco orientado
\begin{equation}
 V_{\rm vac}=
 \begin{pmatrix}21&22\\6&7\end{pmatrix},
 \qquad
 \det V_{\rm vac}=21\cdot7-22\cdot6=15.
 \label{eq:marco-vacancias-21-22-rev11}
\end{equation}
El determinante distinto de cero prueba la independencia de las direcciones
de hueco y retorno. El cardinal \(15\) reaparece así como invariante del
calendario de vacancias. En la ecuación electrónica, la pareja
\((-22,+15)\) procede del selector central demostrado en
\cref{thm:coeficientes-electron-rev6}; el marco
\eqref{eq:marco-vacancias-21-22-rev11} aporta una corroboración aritmética
independiente, con su propia genealogía.

Conviene distinguir dos palabras derivadas del mismo desajuste. La
indicadora
\[
 h_n=\mathbf 1_{\{p_{n+1}-p_n=21\}}
\]
tiene frecuencia \(\sigma_{\rm vac}\) y retornos \(6/7\). La palabra
\[
 z_n=\lceil(n+1)\delta_{\rm vac}\rceil
      -\lceil n\delta_{\rm vac}\rceil
\]
tiene pendiente \(\delta_{\rm vac}\) y contiene \(45\) o \(46\) unidades
en cada bloque de longitud mil. Las leyes \(21/22\), \(6/7\) y \(45/46\)
son lectores distintos de una misma rotación irracional.

\subsection{Polarización acotada y ecuación de continuidad discreta}

Una fase \(\theta\) define
\[
 z_n(\theta)=\lceil(n+1)\delta_{\rm vac}+\theta\rceil
 -\lceil n\delta_{\rm vac}+\theta\rceil.
\]
La polarización intrínseca acumulada es
\begin{equation}
 \begin{aligned}
 P_N(\theta)
 &:=\sum_{n=0}^{N-1}\bigl(z_n(\theta)-\delta_{\rm vac}\bigr)\\
 &=\lceil N\delta_{\rm vac}+\theta\rceil-
   \lceil\theta\rceil-N\delta_{\rm vac},
 \qquad |P_N(\theta)|<1.
 \end{aligned}
 \label{eq:polarizacion-vacancias-rev11}
\end{equation}
La diferencia temporal
\begin{equation}
 J_N(\theta)=P_{N+1}(\theta)-P_N(\theta)
 =z_N(\theta)-\delta_{\rm vac}
 \label{eq:corriente-vacancias-rev11}
\end{equation}
es la corriente aritmética asociada. Las
\cref{eq:polarizacion-vacancias-rev11,eq:corriente-vacancias-rev11}
constituyen una ecuación de continuidad exacta: la carga acumulada permanece
uniformemente acotada y cada pulso es su variación local. La realización
electromagnética aplica a este par una unidad de carga y una escala temporal;
la estructura de polarización y corriente ya está fijada antes de esa carta.

\subsection{Clausura armónica de modos nulos conjugados en el estado electrónico}
\label{sec:nido-armonico-modos-nulos-rev11}
\HMTClaim{MD-R-ELECTRON-NULL-MODE-NEST-CANDIDATE}

La sección electrónica central admite una lectura estructural como clausura
de dos modos nulos conjugados. Para \(u,v>0\), los proyectores hiperbólicos
\(P_\pm\) satisfacen
\[
 N_{\rm sp}(uP_++vP_-)=uv>0,
\]
aunque cada componente aislada tenga norma nula. En esta construcción, la
clausura resonante de modos nulos electromagnéticos conjugados designa el
anillo de transporte de dos hojas cuya holonomía cierra la sección material
central; la emisión y la absorción corresponden a la apertura y al cierre de
una orientación nula respecto de esa sección.

\begin{statusbox}[title={Realización fotónica de la clausura armónica}]
La clausura de dos modos nulos y la positividad de su norma acoplada son
identidades algebraicas. La realización electromagnética se especifica por
un entrelazador con los modos de Maxwell, un Hamiltoniano de clausura y las
tasas de transición correspondientes. La expresión pública
\emph{clausura resonante de modos nulos electromagnéticos conjugados}
nombra el objeto estructural que recibe esa realización.
\end{statusbox}
